% Transcription — fonds Alexandre Grothendieck, shelfmark 68, batch 3 (pages 41-55).
% Established from the facsimile archives/batches/68/batch-03.pdf (a mirror of
% https://grothendieck.umontpellier.fr/68.pdf), per the skill
% transcribe-grothendieck. The batch file carries no cover sheet: PDF page N is
% archivists' page N + 40.
%
% Pass: Opus 5.5 (claude-opus-5-5), 2026-09-22 — first pass, unchecked against the pages by a human.
%
% What the batch is.
%   41-43  End of « Exponentielles (en car. 0) », begun in batch 2 (pp. 35-40):
%          p. 41 closes the proof broken off at the foot of p. 40
%          (delta = rho^N delta, rho nilpotent, so delta = 0); then
%          « Extension aux tenseurs » (e_u (x) e_v = e_{u(x)id + id(x)v}),
%          a « Cor. naïve » (exp u is an algebra automorphism iff u is a
%          derivation), the action on Hom_K(M,N) by delta(f) = vf - fu, the
%          generalisation to K-functors in modules (pp. 42-43), and
%          « Globalisation » to a scheme X of char. 0 (E_Theta).
%          Pencil, faint.
%   44-45  Automorphisms u of a commutative Q-algebra A with (u - id)^{n+1} = 0
%          versus nilpotent derivations: Prop., three corollaries, the
%          bijection L_K(A,I) = G_K(A,I). Pencil; p. 44 carries struck margin
%          conditions (ii'), (iii').
%   46, 48 Blank. Absent.
%   47     « Jeux à N joueurs »: configurations C, players J, j : C -> J,
%          sigma : C -> P(C), initial and terminal configurations, parties,
%          « jeux circulaires ». Pencil.
%   49, 51, 53  Blue ink. Lemmas on sets of cardinal 2; an axiomatic theory
%          of orientation for finite-dimensional real vector spaces
%          (Omega : V -> Ens_2, Stokes maps St_{V'}), uniqueness and
%          existence, reduction to epsilon_n : GL(n,R) -> {+-1},
%          epsilon_n = sg o det (p. 51); p. 53 composition of Stokes maps,
%          products omega x omega', signs (-1)^{dim V . dim V'}, two sketches.
%          The run continues p. 49 -> p. 51 -> p. 53 across the skipped
%          typescripts.
%   50, 52, 54  Three copies of a typescript page headed « XVIII Les
%          épousailles (2) » (paired words: lumière et ombre, …), with
%          handwritten corrections; not mathematics. Absent.
%   55     « Théorie des déterminants gradués »: k commutative ring, category
%          V(k) of free modules of finite type (isomorphisms), graded Picard
%          category P(k), axioms a)-c) for Delta : V -> P, with the
%          commutative square of c). Pencil. Last page of the folder.
%
% Notation fixed for the batch (following batch 2 where it overlaps):
% \mathbb{G}_a and \overline{\mathbb{G}}_a for his G_a and G_a barre;
% \bar e_u, e_u; \mathcal{M} for his script M (a K-functor in modules);
% \mathbb{W}(M) for his outlined W; \mathbf{Aut}, \mathbf{End} for functors
% with outlined initials; \mathbb{E}_\Theta for his outlined E;
% \underline{O}_X as he underlines it; \mathfrak{P}(C) for his gothic P;
% \underline{a} for the underlined involution of a set of cardinal 2;
% \mathcal{V} for his script V (category of vector spaces, then of free
% modules); \mathrm{St} for Stokes maps; \mathbb{P}(k) for his P with a
% doubled stem (graded Picard category).
%
% The dating is the inventory's own for the folder, copied verbatim. Its
% group, [66-89] « Autour de l'enseignement - Thèmes liés (cours et
% directions de recherches) », is dated 1964-1984.
\documentclass[11pt,a4paper]{article}
% Le préambule commun à toutes les transcriptions du fonds Grothendieck.
%
% Il tient en une idée : **l'incertitude se compose, elle ne se résout pas.**
% Une transcription qui lisse un mot douteux a détruit la seule chose pour
% laquelle on la faisait — un lecteur doit pouvoir savoir, à chaque endroit,
% ce qui était sur le feuillet et ce qui a été ajouté. Les six macros qui
% suivent sont donc l'appareil critique, et non de la décoration.
%
% Les mêmes commandes sont reconnues par `scripts/render.mjs`, qui produit la
% vue de lecture du site. Une macro ajoutée ici doit l'être aussi là-bas,
% faute de quoi le rendu échouera bruyamment — ce qui est voulu : un rendu
% silencieusement faux est pire qu'un rendu absent.

\ProvidesPackage{grothendieck}[2026/08/08 Transcriptions du fonds Grothendieck]

\usepackage{amsmath,amssymb,amsthm}
\usepackage{mathrsfs}
\usepackage{stmaryrd}
% Les diagrammes commutatifs sont partout dans ce fonds : c'est la langue
% même de la géométrie algébrique de Grothendieck, et une transcription qui
% les rendrait en prose ne transcrirait rien.
\usepackage{tikz-cd}
% Français par défaut — c'est la langue des feuillets — mais l'anglais est
% chargé aussi : la traduction et le résumé sont des documents anglais, et la
% typographie française (espaces avant les deux-points) y serait une faute.
% Ces documents basculent par \selectlanguage{english} en tête de corps.
\usepackage[english,french]{babel}
\usepackage{fontspec}
\usepackage{geometry}
\geometry{a4paper,margin=25mm,marginparwidth=28mm}
\usepackage{marginnote}
\usepackage{xcolor}
% ulem, not soul. `soul`'s \st and \ul are built on a character-by-character
% scan that XeLaTeX's font machinery breaks: the document fails with an
% undefined control sequence rather than degrading. ulem does the same two jobs
% and is Unicode-engine safe. `normalem` stops it hijacking \emph, which the
% transcriptions use for Grothendieck's own emphasis.
\usepackage[normalem]{ulem}
\usepackage{hyperref}
\hypersetup{colorlinks=true,linkcolor=black,urlcolor=black,pdfborder={0 0 0}}

% --- Métadonnées du lot -----------------------------------------------------
% Reprises telles quelles par le script de rendu, qui les affiche en tête de
% la vue de lecture. Elles fixent ce que le fichier prétend couvrir : sans
% elles, une transcription flotte sans point d'attache dans un fonds de seize
% mille feuillets.

\newcommand{\folder}[1]{\gdef\grothFolder{#1}}
\newcommand{\batch}[1]{\gdef\grothBatch{#1}}
\newcommand{\leaves}[2]{\gdef\grothFirst{#1}\gdef\grothLast{#2}}
\newcommand{\foldertitle}[1]{\gdef\grothFoldertitle{#1}}
\gdef\grothFolder{?}\gdef\grothBatch{?}\gdef\grothFirst{?}\gdef\grothLast{?}\gdef\grothFoldertitle{}

% --- Le résumé --------------------------------------------------------------
%
% Il ouvre la lecture modernisée, sous le titre « Résumé », et il est composé
% en petit. Ce n'est pas un ornement : c'est le seul passage du document qui
% ne soit pas des mathématiques, et un lecteur doit pouvoir le reconnaître
% avant de l'avoir lu — puis le sauter s'il connaît déjà le sujet, ou s'y
% arrêter s'il ne le connaît pas. Le filet qui le suit marque la reprise du
% corps.

\newenvironment{resume}
  {\small\setlength{\parskip}{0.45em}}
  {\par\medskip\hrule\medskip}

% --- Mots-clés ---------------------------------------------------------------
%
% En anglais, en clôture du résumé de la lecture modernisée : le vocabulaire
% moderne sous lequel chercher ce que les feuillets construisent. Le manifeste
% du site (`npm run manifest`) les extrait pour étiqueter la cote entière —
% cette ligne est la source unique des tags, il n'y a pas de fichier à côté.
\newcommand{\keywords}[1]{\par\smallskip\noindent
  {\footnotesize\textsc{Keywords} --- \textit{#1}}\par}

% --- L'appareil critique ----------------------------------------------------

% Le numéro de feuillet, en marge. C'est l'ancre qui permet de régler un
% désaccord sur une formule en regardant *un* feuillet plutôt qu'une chemise
% de six cents. Le site s'en sert aussi pour tourner les pages du fac-similé
% quand on fait défiler la transcription.
\newcommand{\leaf}[1]{%
  \marginnote{\footnotesize\color{gray!70}\textsf{#1}}%
  \label{leaf:#1}\ignorespaces}

% Illisible : on ne devine pas. Un mot inventé qui se lit comme les autres est
% le pire résultat possible.
\newcommand{\ill}{\textcolor{red!60!black}{[\,\ldots\,]}}

% Lecture douteuse : le mot est donné, et signalé comme incertain.
\newcommand{\uncertain}[1]{\uline{#1}}

% Ajout de l'éditeur — une lettre manquante, un mot sous-entendu.
\newcommand{\add}[1]{\textcolor{blue!50!black}{[#1]}}

% Biffé par Grothendieck lui-même. Ce qu'il a rayé fait partie du document :
% une rature dit souvent où la pensée a changé de direction.
\newcommand{\struck}[1]{\sout{#1}}

% Note du transcripteur, sur l'état matériel du feuillet.
\newcommand{\note}[1]{\par\smallskip{\small\color{gray!60!black}\textit{[#1]}}\par\smallskip}

% Note marginale de Grothendieck, distincte du corps du texte.
\newcommand{\marginal}[1]{\par\smallskip\noindent\hspace{1em}%
  {\small\color{gray!60!black}#1}\par\smallskip}

% --- Titre ------------------------------------------------------------------

\AtBeginDocument{%
  \begin{center}
    {\large\bfseries Fonds Alexandre Grothendieck --- cote n\textsuperscript{o}~\grothFolder}\\[2pt]
    {\itshape\grothFoldertitle}\\[4pt]
    {\small feuillets \grothFirst--\grothLast\ (lot \grothBatch)}\\[2pt]
    {\footnotesize\color{gray!60!black}Transcription établie d'après le fac-similé numérisé
     par l'Université de Montpellier.\\
     Les lectures incertaines sont soulignées, les illisibles notées [\,\ldots\,],
     les ajouts entre crochets.}
  \end{center}
  \vspace{1em}}

\endinput


\folder{68}
\batch{3}
\pages{41}{55}
\dating{[à partir de 1978-à partir de 1983]}
\watermark{Édition de démonstration}
\foldertitle{Jeux de position : notes manuscrites (s.d.)}

\begin{document}

\page{41}
\note{la page achève la démonstration interrompue au bas de la p.~40 (lot
précédent) : $\delta = \rho\delta$ avec $\rho$ nilpotent.}

\[
\delta = \rho^N \delta \qquad \forall N \in \mathbb{N}
\]
d'où, prenant $N$ grand \uncertain{pour que} $\rho^N = 0$ :
\[
\delta = 0 \qquad \text{cqfd.}
\]

\section*{Extension aux tenseurs}

$u \in \mathrm{End}_K(M)$, $v \in \mathrm{End}_K(N)$. Alors $M$ ($N$)
$\overline{\mathbb{G}}_a$-Module par $\bar e_u$ ($\bar e_v$), et $M \otimes N$
$\overline{\mathbb{G}}_a$-Module par $\bar e_u \otimes \bar e_v$ — et
\[
\bar e_u \otimes \bar e_v = \bar e_{u \otimes \mathrm{id}_N + \mathrm{id}_M \otimes v}
\]
De même, si $u, v$ nilpotents
\[
e_u \otimes e_v = e_{u \otimes \mathrm{id}_N + \mathrm{id}_M \otimes v}
\]

Extension : $\bigotimes_{i \in I} M_i$ ($I$ ens.\ d'indices fini, de cardinal
$\geq 2$)

\emph{Cor.\ \uncertain{naïve}} Soit $M$ muni d'une structure d'alg.\ (pas
néc.\ ass.\ ni unitaire, encore moins commutative), sur $K$,
$u \in \mathrm{End}_{K\text{-mod}}(M)$, donc $\overline{\mathbb{G}}_a$ (resp.\
$\mathbb{G}_a$ si $u$ nilpotent) opère sur le $K$-module $M$. Pour que cette
opération soit une op.\ d'automorph.\ d'algèbre, il f.\ et s.\ que $u$ soit une
\emph{dérivation}.

\page{42}
Si $u$ nilpotent, il f.\ et s.\ que $\exp u$ soit un automorph.\ d'algèbre.

Op.\ sur $\mathrm{Hom}_K(M, N)$, quand $\overline{\mathbb{G}}_a$ (resp.\
$\mathbb{G}_a$) opère sur $M$ par $\bar e_u$ (resp.\ $e_u$), sur $N$ par
$\bar e_v$ (resp.\ $e_v$). Alors il opère sur $\mathrm{Hom}_K(M, N)$ par
$\bar e_\delta$ ($e_\delta$), où
\[
\delta(f) = vf - fu
\]

Mais en tt rigueur $\bar e_\delta$ et $e_\delta$ ne sont définis \ill{} que si
$\mathrm{Hom}_K(M, N)$ est \uncertain{is.\ à} un foncteur de la forme
$\mathbb{W}(M)$ — ce qui n'est guère assuré que si $M$ proj.\ de t.f.
\note{« is.\ à » est un ajout interlinéaire, marqué d'un signe d'insertion.}

Mais on peut généraliser théorie de l'exponentielle \struck{aux tenseurs, les}
aux endomorphismes de tt $K$-foncteurs $\mathcal{M}$ en
\uncertain{$\mathbb{Q}$}-modules
\note{la lettre est surchargée et entourée ; on attend peut-être $K$.}
\[
\mathcal{M}(K') \text{ est un } K'\text{-module}
\]
\[
\mathcal{M}(K') \to \mathcal{M}(K'') \quad (K' \to K'')\text{-semi-linéaire}
\]

Alors $\mathrm{End}_K(\mathcal{M})$ est une $K$-\emph{algèbre}.

a) Si $u \in N(\mathrm{End}_K(\mathcal{M}))$, $\exp u \in
\mathrm{End}_K(\mathcal{M})$ et aussi $\exp(\lambda u)$
$\forall \lambda \in \mathbb{G}_a(K') \simeq K'$ défini,
$\in \mathrm{End}_{K'}(\mathcal{M}_{K'})$. On trouve sur $\mathcal{M}$ une
structure de $\mathbb{G}_a$-Module \ill{}

\page{43}
b) Si $u \in \mathrm{End}_K(\mathcal{M})$, alors on définit
\[
\bar e_u : \overline{\mathbb{G}}_{a,K} \to \mathbf{Aut}_K(\mathcal{M})
\]
\note{la flèche aboutit à un symbole surchargé, à initiale doublée, lu
$\mathbf{Aut}$.}
\[
\bar e_u(\lambda) = \exp(\lambda u)
\]

\emph{NB} La question se pose encore si $u \mapsto \bar e_u$ et
\marginal{?}
$u \mapsto e_u$ sont des \emph{bij.} :
\[
\mathrm{End}_K(\mathcal{M}) \to \mathrm{Hom\,gr}(\overline{\mathbb{G}}_{a,K}, \mathbf{Aut}_K(\mathcal{M}))
\]
\[
N(\mathrm{End}_K(\mathcal{M})) \to \mathrm{Hom\,gr}(\mathbb{G}_{a,K}, \mathbf{Aut}_K(\mathcal{M}))
\]

Extension aux \struck{\ill{}} algèbres symétriques, tensorielles,
extérieures …

\section*{Globalisation}

\struck{S} $X$ schéma de car.\ 0, $\Theta$ \struck{S} dérivation de
$\underline{O}_X$, supposons $\Theta$ loc.\ nilpotente. Alors $\exp \Theta$ est
un automorphisme de $\underline{O}_X$ qui induit l'identité sur
$\underline{O}_{X_{\mathrm{red}}} = \underline{O}_X / N\underline{O}_X$, et
correspond à un automorphisme de $X$, que nous désignons par
$\mathbb{E}_\Theta$. Si $\Theta$ est une $S$-dérivation, $\mathbb{E}_\Theta$
est un $S$-automorphisme. Enfin
\[
\begin{cases}
\mathbb{E}_0 = \mathrm{id} \\
\mathbb{E}_{\Theta + \Theta'} = \mathbb{E}_\Theta + \mathbb{E}_{\Theta'} & \text{si } [\Theta, \Theta'] = 0
\end{cases}
\]
\note{tel qu'écrit : le signe entre $\mathbb{E}_\Theta$ et $\mathbb{E}_{\Theta'}$
a la forme d'un $+$ ; pour des automorphismes on attend une composition.}

\page{44}
\emph{Prop} $A$ $\mathbb{Q}$-alg.\ comm.\ \uncertain{unit.},
$u$ un automorph.\ de $\mathbb{Q}$-alg., $n \in \mathbb{N}$, $J$ l'idéal
engendré par les $u(x) - x$ ($x \in A$)
\marginal{\struck{NB $u$ induit l'identité sur $\mathrm{gr}_J(A)$}}
\note{la note marginale en haut à droite est barrée d'un trait oblique.}

1) Conditions équivalentes
\begin{itemize}
  \item[(i)] $u$ op.\ \uncertain{diff.}\ de $A$ d'ordre $\leq n$
  \item[(ii)] $J^{n+1} = 0$
  \item[(iii)] $\exists\, I$ idéal de $A$ tel que $I^{n+1} = 0$, et que $u$
  induise l'identité dans $A/I$, i.e.\ $ux - x \in I$ $\forall x \in A$
  [i.e.\ $I \supset J$]
  \item[\struck{(iv)}] $\exists$ dérivation $\Theta : A \to A$ avec
  $\Theta^{n+1} = 0$, $u = \exp \Theta$
  \marginal{est-il nécessaire ?}
\end{itemize}
\note{au-dessus de la ligne (iv), en interligne : « Pour ceci, il suffit que
$(u - \mathrm{id})^{n+1} = 0$, ou (ce qui revient au même) ».}

\marginal{\struck{(ii') $u$ induit l'identité} \struck{(iii') $\exists\, I$
\ill{} que $I^{n+1} = 0$, et que $u$ induise l'identité \ill{}
$\mathrm{gr}_I(A)$}}
\note{les deux conditions de la marge gauche sont barrées et encadrées.}

\struck{NB Si $N(A)$ est nilpotent, ces conditions équivalentes}

2) L'application $\Theta \mapsto \exp \Theta$ est une bijection de l'ens.\ des
dérivations $\Theta$ de $A$ telles que $\Theta^{n+1} = 0$, dans l'ens.\ des
automorphismes $u$ de $A$ telles que $(u - \mathrm{id})^{n+1} = 0$ (\ill{} de
$A$ (\ill{}) \ill{} plus \ill{}).

\emph{Cor} Conditions équivalentes sur un automorph.\ $u$ de $A$ :
\begin{itemize}
  \item[(i)] $u$ op.\ \uncertain{diff.}\ de $A$
  \item[(ii)] $J$ idéal nilp.
  \item[(iii)] $\exists\, I$ nilp.\ de $A$ tel que $u$ induise
  $\mathrm{id}_{A/I}$
  \item[(iv)] $\exists\, \Theta$ dérivation nilp.\ de $A$, telle que
  $u = \exp \Theta$
\end{itemize}
De plus, $\Theta \mapsto \exp \Theta$ est une bij.\ de l'ens.\ des dérivations
nilpotentes de $A$ et l'ens.\ des autom.\ de $A$ qui sont des op.\
\uncertain{diff.}\ de $A$.

\page{45}
\emph{Cor} \struck{A} Supposons $NA$ nilp.\ \ill{} (p.\ ex.\ $A$ noeth.).
Alors les conditions précédentes sur $u$ équivalent à :
\begin{itemize}
  \item[(v)] $u$ induit l'identité sur $A_{\mathrm{red}}$
\end{itemize}
\uncertain{Sous} ces conditions, et si $N(A)^{N+1} = 0$, les conditions de la
prop sont satisfaites pour $(u, N)$, i.e.\ $u$ \uncertain{opérateur diff.}\ de
$A$ est d'ordre $\leq N$.

\emph{Cor.} Supposons $A$ une $K$-alg ($K$ une $\mathbb{Q}$-alg). Si
$u = \exp \Theta$ ($\Theta$ dérivation nilpotente) \uncertain{alors} $u$ est un
$K$-automorph.\ ssi $\Theta$ est une $K$-dérivation.

\emph{Cor} Soit $I$ un idéal nilpotent de \struck{A} \ill{} $K$-alg.\ comm.\
$A$ ($K$ $\mathbb{Q}$-alg.\ commut., $I^{N+1} = 0$, \ill{}), soit
$L_K(A, I)$ l'ens.\ des $K$-dérivations \struck{$\Theta$} de $A$ qui induisent
0 dans $A/I$ (i.e.\ $\Theta(A) \subset I$), $G_K(A, I)$ le gr.\ des
$K$-automorph.\ de $A$ \struck{telles} qui induisent l'identité sur $A/I$ (i.e.\
$(u - \mathrm{id})(A) \subset I$ i.e.\ $u(x) - x \in I$ $\forall x \in A$).
Alors
\begin{itemize}
  \item[a)] $\forall \Theta \in L_K(A, I)$, on a $\Theta^{N+1} = 0$ (plus
  génér.\ ment, $\forall i \in \mathbb{N}$, $\Theta^i A \subset I^i$)
  \item[b)] $\Theta \mapsto \exp \Theta : L_K(A, I) \simeq G_K(A, I)$
\end{itemize}

\page{47}
\section*{Jeux à $N$ joueurs}

\begin{itemize}
  \item[] $C$ ensemble des configurations (ou positions)
  \item[] $J$ ensemble des joueurs ($\mathrm{Card}\, J = N$)
  \item[] $j : C \to J$ application (pour toute configuration, il y a un
  joueur « dont c'est le tour » \ill{})
  \item[] $\sigma : C \to \mathfrak{P}(C)$ pour toute configuration $x$, il y a
  l'ensemble $\sigma(x)$ des configurations « suivantes »
  \item[] [$\Sigma \subset C \times C$ une relation \quad [N.B.\
  $(x, y) \in \Sigma \iff y \in \sigma(x)$]]
  \item[] $x_0 \in C$ configuration initiale.
  \item[] $C_0 \subset C$ ens.\ des configurations \struck{gagnantes}
  \emph{terminales} i.e.\ $C_0 = \{ x \in C \mid \sigma(x) = \emptyset \}$
\end{itemize}
\note{la lettre $j$ de l'application est surchargée.}

\struck{Jeu} \emph{Partie} soit $j_0 = j(x_0)$ (joueur initial)
\begin{itemize}
  \item[1)] $j_0$ « choisit » un $x_1 \in \sigma(x_0)$. Soit $j_1 = j(x_1)$
  \item[2)] $j_1$ « choisit » un $x_2 \in \sigma(x_1)$ soit $j_2 = j(x_2)$
  \item[—]
  \item[$n$)] $j_{n-1}$ « choisit » $x_n \in \sigma(x_{n-1})$
\end{itemize}

Partie à $n$ coups : \struck{\ill{}} $(x_1, \ldots, x_n) \in C^n$, \quad
$x_i \in \sigma(x_{i-1})$ \quad $1 \leq i \leq n$

\emph{Jeux circulaires}
\begin{itemize}
  \item[1)] $j$ surjectif (tout joueur est effectif)
  \item[2)] $\exists\, \rho : J \to J$ telle que :
  \struck{$\forall x \in C$, $y \in \sigma(x)$}
  $\forall (x, y) \in \Sigma \subset C \times C$, … $j(y) = \rho(j(x))$
\end{itemize}
\emph{NB} $\rho$ unique si on suppose que $\forall x \in C$,
$\sigma(x) \neq \emptyset$, ou plus génér.\ que $\forall j \in J$,
$\exists\, x \in C$ tel que $\sigma(x) \neq \emptyset$, $j(x) = j$
\begin{itemize}
  \item[3)] $\rho$ est une permutation circ.
\end{itemize}

\page{49}
\note{encre bleue, jusqu'à la p.~53.}

\emph{Lemme 1} $P, Q$ ens.\ de card.\ 2, ($x \in P$ \struck{\ill{}}
$x' = \underline{a}_P x$), $\varphi : P \to Q$ conditions équivalentes
\begin{itemize}
  \item[a)] $\varphi$ bij \quad b) $\varphi$ inj \quad c) $\varphi$ surj
  \quad d) $\varphi$ non-constante
  \item[e)] $\varphi$ commute à $\mathbb{Z}_2$ \quad f)
  $\varphi(x') = \underline{a}_Q \varphi(x)$
\end{itemize}

\emph{Lemme 2} $P, Q, R$ ens.\ de card.\ 2, $\varphi : P \times Q \to R$
\uncertain{appl.}, $x \in P$, $x' = \underline{a}\, x$

Conditions équivalentes
\begin{itemize}
  \item[a)] $\varphi$ se factorise en iso $P \wedge Q \to R$
  \item[b)] $\varphi_x : y \mapsto \varphi(x, y) : Q \to R$ est iso, et
  $\varphi_{x'}$ \struck{$\varphi_x$} $= \underline{a}_R \varphi_x$
\end{itemize}

Théorie orientation pour \struck{\ill{}} \uncertain{vectoriels} réels de dim.\
finie (cat.\ $\mathcal{V}$, avec \emph{iso})
\begin{itemize}
  \item[a)] \struck{\ill{}} $\Omega : \mathcal{V} \to (\mathrm{Ens}_2)$
  \item[b)] $\Omega(0) \simeq \{\pm 1\}$
  \item[c)] $\forall V \in \mathcal{V}$, \struck{\ill{}} $W \subset V$
  hyperplan, $V'$ demi-espace de $V$ de frontière $W$, on a donc une
  application de Stokes
  \[
  \mathrm{St}_{V'} : \Omega(V) \to \Omega(W) \qquad \text{fonctorielle en } (V, V')
  \]
\end{itemize}

\emph{Axiomes}
\begin{itemize}
  \item[1)] $\mathrm{St}_{V'}$ bijectif \struck{(cf lemme 1)}
  \uncertain{car} $\forall V \in \mathcal{V}$, $\mathrm{Card}\, \Omega(V) = 2$
  \item[2)] Si $V''$ est l'opposé de $V'$, …
  $\mathrm{St}_{V''} = \underline{a} \circ \mathrm{St}_V$
\end{itemize}
\note{tel qu'écrit : $\mathrm{St}_V$ ; on attend $\mathrm{St}_{V'}$.}

\emph{Axiomes équivalents}
\begin{itemize}
  \item[1)] $\forall V \in \mathcal{V}$, $\mathrm{card}\, \Omega(V) = 2$
  \item[2)] $\forall V \in \mathcal{V}$, $W \in \mathrm{Hyp}(V)$,
  \struck{2)} l'application
  \struck{$\Omega(W, V)$} $\mathrm{Dem}_{W,V} \times \Omega(V) \to \Omega(W)$,
  $(V', \omega) \mapsto \mathrm{St}_{V'}(\omega)$ se factorise en une
  \struck{application} iso
  \[
  \Omega(W, V) \wedge \Omega(V) \xrightarrow{\sim} \Omega(W)
  \]
  $\to$ fonctorielle en $(V, W)$
\end{itemize}
\note{l'indice de « $\mathrm{Dem}$ » est lu sous une rature ; la lecture
$\mathrm{Dem}_{W,V}$ (demi-espaces) est incertaine.}

\struck{Proposition} Hom d'une théorie $(\Omega, \mathrm{St})$ dans théorie
$(\Omega', \mathrm{St}')$ : homs de foncteurs $\varphi : \Omega \to \Omega'$
compatible avec b), et c)

\emph{Th.}
\begin{itemize}
  \item[a)] Tout hom d'une théorie dans une autre est un iso.
  \item[b)] Pour 2 théories de l'orientation, $\exists !$ isom de l'une dans
  l'autre
  \item[c)] $\exists$ théorie de l'orientation.
\end{itemize}

\emph{Dém}
\begin{itemize}
  \item[a)] $\varphi_V : \Omega(V) \simeq \Omega'(V)$ iso, par réc.\ sur dim.,
  utilisant Stokes
  \item[b)] De $=$, unicité de $\varphi$ par récurrence sur dim. Existence
  \ill{} \struck{Détermination d'une thé} Existence d'une théorie de
  l'orientation ?
\end{itemize}

donner $\Omega$ $\iff$ donner $(\Omega_n)_{n \in \mathbb{N}}$,
$\Omega_n : \mathcal{V}_n \to \mathrm{Ens}_2\ \mathrm{Tors}(\mathbb{Z}/2)$
\note{au-dessus de la flèche, en interligne : « $\mathcal{V} \to
\mathrm{Ens}_2$ ».}

donner $\Omega_n$ $\iff$ donner
$\varepsilon_n : GL(n, \mathbb{R}) = G_n \to \mathbb{Z}/2$, \emph{et}
$\omega_n = \Omega(\mathbb{R}^n)$

donner $\mathrm{St}$ $\iff$ $(\mathrm{St}_n)_{n \in \mathbb{N}^*}$, où
$\mathrm{St}_n$ est donnée de Stokes en dim.\ $n$
\[
\Omega(V, W) \wedge \Omega(W) \xrightarrow{\sim} \Omega(W) \qquad \text{pour } \dim V = n
\]
\note{tel qu'écrit, le second facteur semble être $\Omega(W)$ ; on attend
$\Omega(V)$.}

donner \struck{$\mathrm{St}_n$} $(\mathrm{St}_n)_{n}$ revient à donner

\page{51}
\note{suite de la p.~49 ; la p.~50 est un tapuscrit sans rapport.}

\[
\underbrace{\Omega(\mathbb{R}^{n-1}, \mathbb{R}^n)}_{\mathbb{1}} \wedge \Omega(\mathbb{R}^n) \to \Omega(\mathbb{R}^{n-1}) \qquad \text{i.e.} \qquad \omega_n \xrightarrow{\sim} \omega_{n-1}
\]
\marginal{\uncertain{choisi} $\mathbb{R}^{n-1} \times \mathbb{R}^+$ comme
$\tfrac{1}{2}$ espace de $\mathbb{R}^n$}

compatible avec $\mathrm{Aut}(\mathbb{R}^{n-1}, \mathbb{R}^n) =$ ss-gr.\ de
$GL(n, \mathbb{R})$ des matrices
\[
u_n = \begin{pmatrix} u_{n-1} & \vdots \\ 0 \ \cdots \ 0 & \lambda \end{pmatrix}
\]
opérant sur $\omega_{n-1}$ par \struck{$\varepsilon(u_{n-1})$}, sur $\omega_n$
par \struck{\ill{}} $\varepsilon_n(u_n)$. $\mathrm{sg}\,\lambda =$
\struck{$\mathrm{sg}\det$ \ill{}} condition qui signifie des conditions
\[
\varepsilon_n(u_n) = \mathrm{sg}(\lambda)\, \varepsilon_{n-1}(u_{n-1})
\]

\struck{qui est donc \ill{} satisfaisante ! Donc la donnée d'une théorie de
l'orient.\ équivaut à la donnée de $\{(\varepsilon_n)_{n \in \mathbb{N}}$ et
$\varepsilon_n : G_n \to \{\pm 1\}$ \ill{} $s : \omega_n \xrightarrow{\sim}
\omega_{n-1}$, les $\omega_n \simeq \omega_{n-1}$, et
$\omega_0 \simeq \{\pm 1\}$, donc $\omega_n \simeq \{\pm 1\}\}$}
\note{tout ce passage est encadré et biffé d'une grande croix.}

Ceci détermine $\varepsilon_1 : \mathbb{R}^* \to \{\pm 1\}$ i.e.
\[
\varepsilon_1 = \mathrm{sg} : \mathbb{R}^* \to \{\pm 1\}
\]
et de proche en proche, par récurrence
\[
\varepsilon_n(u_n) = \mathrm{sg} \circ \det(u_n) \qquad u \in GL(n, \mathbb{R})
\]

\struck{Lemme 2}

\page{53}
\[
U \underbrace{\subset}_{\omega'} W \underbrace{\subset}_{\omega} V \qquad \Omega(V/W) \wedge \Omega(V) \to \Omega(W)
\]
\[
\mathrm{St}_\omega : \Omega(V) \to \Omega(W)
\]
\[
\mathrm{St}_{\omega'} : \Omega(W) \to \Omega(U)
\]
\[
\mathrm{St}_{\omega'} \mathrm{St}_\omega : \Omega(V) \to \Omega(U)
\]
\[
I(\omega, \omega') \in \Omega
\]

\marginal{$(\omega, \omega') \mapsto \omega \times \omega'$}
\[
\Omega(V) \wedge \Omega(V') \to \Omega(V \times V')
\]
\[
\omega' \times \omega = (-1)^{\dim V \cdot \dim V'}\, \omega \times \omega'
\]
\note{un symbole est raturé devant $\omega \times \omega'$.}

\struck{$\mathrm{St}(\omega)\ \mathrm{St}(\omega')$}
\[
\begin{array}{ccccccccc}
\mathbb{R} & \subset & \mathbb{R}^2 & \subset & \mathbb{R}^3 & \text{—} & \mathbb{R}^{n-1} & \subset & \mathbb{R}^n \\
\cup & & \cup & & \cup & & \cup & & \cup \\
\mathbb{R}^+ & \subset & \mathbb{R} \times \mathbb{R}^+ & \subset & \mathbb{R}^2 \times \mathbb{R}^+ & \subset & \mathbb{R}^{n-2} \times \mathbb{R}^+ & \subset & \mathbb{R}^{n-1} \times \mathbb{R}^+
\end{array}
\]

\struck{$e_1 \wedge \cdots$} $e_{n-1} \wedge e_n$ \qquad
$e_1 \wedge \cdots \wedge e_{n-1} \wedge e_n$

$\omega \times \omega'$ \qquad $I(\omega, \omega')$ \qquad $I(\omega', \omega)$

\[
\begin{cases}
\mathrm{St}(\omega \times \omega') = \mathrm{St}(\omega) \times \omega' \\
\mathrm{St}(\omega \times \omega') = \omega \times \mathrm{St}\,\omega'
\end{cases}
\]
\note{au-dessus de la première ligne, relié par un trait : « $(-1)^{?}$ » ; dans la seconde, un symbole est raturé devant $\omega \times \mathrm{St}\,\omega'$.}

$e_1\, e_2 \ldots e_n$

$\underbrace{e_1, \ldots, e_p}\ f_1 \ldots f_q$
\[
\mathrm{St}(e_1 \ldots e_{p-1}\, e_p\, f_1 \ldots f_q) = (-1)^q\, e_1 \ldots e_{p-1}\, \hat e_p\, f_1 \ldots f_q
\]
\note{au bas de la page, deux croquis : un carré, et un quadrant dont les
deux demi-axes portent des flèches et dont l'intérieur est hachuré,
accompagné d'une flèche courbe de rotation.}

\page{55}
$k$ anneau commut.

$\mathcal{V} = \mathcal{V}(k)$ catégorie des Modules libres de type fini sur
$k$ (morphismes : \emph{iso})
\[
\mathcal{V} = \coprod_{n \in \mathbb{N}} \mathcal{V}_n(k)
\]
$\mathcal{V}_n(k)$ catégorie des modules libres de rang $n$ sur $k$ \quad
$\mathcal{V}$ catégorie additive (graduée)

\struck{$\mathcal{V} \xrightarrow{\Delta} \mathcal{V}_{\mathbb{Z}}$ foncteur}

\struck{\ill{}} $\mathbb{P}(k)$ catégorie de Picard des $k$-modules libres de
rang 1 $\mathbb{Z}$-gradués, avec donnée de commut.\ l'isomorphisme de
Koszul

\section*{Théorie des déterminants gradués}

\note{titre souligné ; un premier « déterminants » y est barré, et un mot lu \uncertain{modules} est écrit au-dessus.}

\begin{itemize}
  \item[a)] $\mathcal{V} \xrightarrow{\Delta} \mathbb{P}$ \quad
  $\otimes$-foncteur [donnée d'iso fonct.
  \[
  \Delta(V + V') \overset{\varphi_{V,V'}}{\simeq} \Delta(V) \otimes \Delta(V')
  \]
  ($\Rightarrow \Delta(0) \simeq \mathbb{1}$ can.)
  \item[b)] \struck{\ill{}} Si $\mathrm{rg}(V) = 1$, $\Delta(V) \simeq V$, de
  degré 1
  \item[c)] Soient $V \supset W_1$ ; $W_2$, $W_2'$ deux suppl.\ de $W_1$ dans
  $V$ ($W_1$, $W_2$ libres dans $V$, $W_2'$ …, \uncertain{ni})
  \[
  \alpha : W_1 + W_2 \xrightarrow{\sim} V \qquad
  \alpha' : W_1 + W_2' \xrightarrow{\sim} V \qquad
  \lambda : W_2 \simeq W_2'
  \]
\end{itemize}

\begin{tikzcd}
\Delta(W_1 + W_2) \arrow[r, "\Delta(\alpha)"] \arrow[d, "\varphi_{W_1, W_2}"'] & \Delta(V) & \Delta(W_1 + W_2') \arrow[l, "\Delta(\alpha')"'] \arrow[d, "\varphi_{W_1, W_2'}"] \\
\Delta(W_1) \otimes \Delta(W_2) & & \Delta(W_1) \otimes \Delta(W_2') \arrow[ll, "\mathrm{id} \otimes \Delta(\lambda)"]
\end{tikzcd}

\note{les quatre flèches du carré portent en outre le signe $\simeq$ ; à
gauche, l'indice est écrit $\varphi_{W_1, W_1}$, lapsus pour
$\varphi_{W_1, W_2}$.}

commutatif.]

[Équivaut à : si $u$ est un automorph.\ de $V$ qui est l'identité sur $W_1$ et
sur $V/W_1$ ($\simeq W_2$), alors $\Delta(u) = \mathrm{id}_{\Delta(V)}$]

\end{document}
